\documentclass[11pt,a4paper]{article}
\usepackage[margin=2.2cm]{geometry}
\usepackage[T1]{fontenc}
\usepackage{lmodern}
\usepackage{amsmath,amssymb}
\usepackage{booktabs,longtable,array}
\usepackage{xcolor}
\usepackage{hyperref}
\usepackage{enumitem}
\usepackage{microtype}

\definecolor{riskblue}{HTML}{1976A3}
\definecolor{riskred}{HTML}{B83A3A}
\hypersetup{colorlinks=true,linkcolor=riskblue,urlcolor=riskblue}
\setlist{nosep}
\newcommand{\RiskZero}{\textsc{Risk0}}
\newcommand{\Prob}{\mathbb{P}}
\newcommand{\E}{\mathbb{E}}
\newcommand{\ind}{\mathbf{1}}

\title{\RiskZero{}: Algorithms and Computational Formulas}
\author{Morpho Risk Radar MVP}
\date{10 October 2026}

\begin{document}
\maketitle

\begin{abstract}
This document specifies the algorithms used by the \RiskZero{} MVP to discover
Morpho dependencies, compare vaults, evaluate entity health, propagate incidents,
estimate 30-day scenario risk, estimate immediate exit cost, and produce an
explainable \texttt{STAY}, \texttt{EXIT}, or \texttt{UNAVAILABLE} result.  It
distinguishes empirical probabilities from transparent model estimates and makes
missing data an explicit state rather than silently replacing it by zero.
\end{abstract}

\tableofcontents

\section{Scope, notation, and invariants}

The monitored dependency system is a directed attributed graph
\begin{equation}
G_t=(V_t,E_t,X_t),
\end{equation}
where $V_t$ contains chain, protocol, vault, market, asset, oracle, interest-rate
model (IRM), and manually entered position entities; $E_t$ contains typed
dependencies; and $X_t$ contains time-stamped observations and health state.
Relevant edge types include \texttt{allocated\_to}, \texttt{lends},
\texttt{borrows}, \texttt{uses\_collateral}, \texttt{uses\_oracle},
\texttt{uses\_irm}, \texttt{deployed\_on}, and \texttt{similar\_to}.

Every observation stores an observation time $t_o$, retrieval time $t_r$, source
identity, freshness class, metrics, and an explicit missing-field set.  Three
invariants govern all calculations:
\begin{enumerate}
  \item unknown critical data is never interpreted as zero risk;
  \item empirical and model-estimated probabilities are labelled separately;
  \item every incident impact retains at least one explainable dependency path.
\end{enumerate}

The decision horizon is $H=30$ days.  Decimal values cross the service boundary
as decimal strings to avoid binary floating-point ambiguity.

\section{Collection and scheduling}

Two non-overlapping loops update the radar:
\begin{align}
\Delta_{\mathrm{fast}} &= 300\ \mathrm{s},\\
\Delta_{\mathrm{slow}} &= 3600\ \mathrm{s}.
\end{align}
The fast loop collects changing market, oracle, vault, and chain state, evaluates
health, updates incidents, and propagates impacts.  The slow loop rediscovers
topology and recomputes vault similarity.  For loop $\ell$, an asynchronous lock
$L_\ell$ enforces
\begin{equation}
\ind[\text{start next run of }\ell]=\ind[L_\ell\text{ is free}].
\end{equation}
If the lock is held, the tick is stored as \texttt{skipped}; overlapping work is
not queued.  Failures are isolated into a recorded scheduler run and do not
terminate the other loop.

Raw snapshots and superseded health records have 30-day retention.  Hourly
aggregates have 365-day retention.  Current health and incident history survive
pruning.  This supports a local SQLite MVP; a distributed deployment would need
a shared scheduler lock, message broker, durable event log, and managed backups.

\section{Morpho market representation}

A Morpho market $m$ is represented by the tuple
\begin{equation}
m=(a_L,a_C,o,r,\lambda),
\end{equation}
where $a_L$ is the loan asset, $a_C$ the collateral asset, $o$ the oracle, $r$
the IRM, and $\lambda\in(0,1)$ the liquidation loan-to-value (LLTV).  Vault
allocations add weighted edges
\begin{equation}
v \xrightarrow[\ w_{vm}\ ]{\texttt{allocated\_to}} m,
\qquad 0\leq w_{vm}\leq 1,
\end{equation}
with $w_{vm}$ interpreted as the observable allocation share when available.
The graph then links $m$ to its loan asset, collateral asset, oracle, IRM,
protocol, and deployment chain.

\section{Vault similarity}

Similarity is computed between a catalog vault $c$ and candidate vault $v$.
Let $J_w(A,B)$ be weighted Jaccard similarity,
\begin{equation}
J_w(A,B)=\frac{\sum_{k\in A\cup B}\min(a_k,b_k)}
                   {\sum_{k\in A\cup B}\max(a_k,b_k)},
\end{equation}
with zero contribution when the denominator is zero.  The MVP score is
\begin{align}
S(c,v)={}&0.35J_w(M_c,M_v)+0.20J(A_c,A_v)+0.15J(O_c,O_v)\nonumber\\
        &+0.10J(R_c,R_v)+0.10P_{\lambda}(c,v)
          +0.05I_{\mathrm{chain}}+0.05I_{\mathrm{curator}},
\end{align}
where $M,A,O,R$ denote market, asset, oracle, and IRM sets.  The LLTV proximity
component is
\begin{equation}
P_{\lambda}(c,v)=\max\left(0,1-|\bar\lambda_c-\bar\lambda_v|\right).
\end{equation}
Missing dimensions are recorded and contribute zero; they are not imputed.

Candidates are stored when $S>0.33$.  A maximum of ten candidates per catalog
vault are normally visible when $S>0.60$.  Stored candidates below the visual
threshold remain processed and are revealed if an incident path reaches them.

Allocation concentration is summarized with the Herfindahl--Hirschman index
\begin{equation}
\mathrm{HHI}(v)=\sum_m w_{vm}^{2}.
\end{equation}
The current policy flags concentration when $\mathrm{HHI}>0.25$.

\section{Oracle consensus and deviation}

For oracle $o_i$ representing a common economic pair at time $t$, normalized
prices are compared with the peer median
\begin{equation}
\widetilde p_t=\operatorname{median}\{p_{1,t},\ldots,p_{n,t}\}.
\end{equation}
The relative deviation of oracle $i$ is
\begin{equation}
d_{i,t}=\frac{|p_{i,t}-\widetilde p_t|}{\widetilde p_t}.
\end{equation}
The MVP policy is
\begin{equation}
h_o(d)=
\begin{cases}
\texttt{CRITICAL}, & d\geq 0.05,\\
\texttt{WARNING}, & 0.02\leq d<0.05,\\
\texttt{HEALTHY}, & d<0.02,\\
\texttt{UNKNOWN}, & \text{insufficient, stale, or incompatible references.}
\end{cases}
\end{equation}
Staleness, invalid rounds, and unavailable peer consensus are evaluated in
parallel with price deviation.  The most severe valid condition wins.

\section{Entity health state machine}

Each rule emits a proposed state $\hat h_{e,t}$, score, reasons, missing fields,
and a version.  The externally visible state is
\begin{equation}
h_{e,t}\in\{\texttt{HEALTHY},\texttt{WARNING},\texttt{CRITICAL},\texttt{UNKNOWN}\}.
\end{equation}
\texttt{UNKNOWN} is not healthy and is visually distinct.  To reduce flapping,
a warning requires two consecutive violating cycles and recovery requires three
consecutive healthy cycles.

For a market with supplied amount $S$ and borrowed amount $B$, utilization is
\begin{equation}
u=\begin{cases}B/S,&S>0,\\\text{unknown},&S\leq0.\end{cases}
\end{equation}
The current utilization thresholds are warning at $u\geq0.90$ and critical at
$u\geq0.98$.  Bad debt is critical.  Near-liquidation borrower ratios use
warning and critical portfolio-share thresholds of $0.10$ and $0.30$.

For vault allocations, the unmodelled share is
\begin{equation}
q_{\mathrm{unknown}}=\max\left(0,1-\sum_{m\in M_{\mathrm{modelled}}}w_{vm}\right).
\end{equation}
When $q_{\mathrm{unknown}}>0.05$, health cannot be asserted without qualification;
the missing allocation coverage is carried into the state and recommendation.

\section{Incident lifecycle}

A transition to warning or critical opens or updates an incident keyed by the
source entity and rule.  Duplicate evaluations update the same incident rather
than producing notification storms.  A recovery sequence resolves the incident
but does not delete its impacts or notification history.  The persistence rule
is therefore
\begin{equation}
\text{resolved}\neq\text{forgotten}.
\end{equation}
Notifications are durable local records with optional read time.  Toasts are a
transient projection of those records and may disappear without deleting them.

\section{Reverse dependency propagation}

An incident originating at $s$ propagates against dependency direction: if a
vault depends on a market and that market depends on an oracle, an oracle
incident travels oracle $\rightarrow$ market $\rightarrow$ vault.  The MVP uses
breadth-first search over the reverse adjacency relation
\begin{equation}
\operatorname{pred}(x)=\{u:(u,x)\in E\}.
\end{equation}
For each reached target $y$, BFS stores a shortest path
\begin{equation}
\pi(s,y)=(s,\ldots,y),\qquad d(s,y)=|\pi|-1.
\end{equation}
Alternative shortest paths are counted for explainability.  Distance confidence
is deliberately simple and bounded:
\begin{equation}
C(s,y)=\max(0.50,1-0.05d(s,y)).
\end{equation}

For lender exposure amount $A$ attached to a vault, an affected amount along an
allocation path $\pi$ is
\begin{equation}
I_{\mathrm{lender}}(\pi)=A\prod_{e\in\pi}w_e,
\end{equation}
where unweighted dependency edges contribute one.  For a borrower, debt amount
is the primary at-risk amount while collateral and LTV drive liquidation tests.
If several paths reach the same target, amounts are not blindly summed without
deduplication; one primary path and an alternate-path count are retained.

\section{Scenario impact}

For scenario $j$ with severity $s_j$, deterministic propagation produces gross
loss $L_j$ and a hard-trigger set $T_j$.  A lender loss on an affected allocation
may be approximated by
\begin{equation}
L_j=A\sum_{m\in M_j}w_{vm}\ell_{jm},
\end{equation}
where $\ell_{jm}\in[0,1]$ is the scenario loss fraction for market $m$.

For a borrower with collateral value $C$, debt value $D$, current LTV
$\rho=D/C$, and a collateral shock $s\in[0,1)$,
\begin{equation}
\rho'=\frac{D}{C(1-s)}=\frac{\rho}{1-s}.
\end{equation}
Liquidation is triggered when $\rho'\geq\lambda$.  Let $C_o$ be collateral
valued at the oracle price after the shock and let $f>1$ be the liquidation
incentive factor.  The engine first computes collateral seized at oracle value,
the corresponding repaid debt, and collateral seized at the stressed value:
\begin{equation}
q_o=\min(C_o,D'f),\qquad R=\frac{q_o}{f},
\qquad
Q=\begin{cases}
q_o, & C_o=C',\\
q_o\,C'/C_o, & C_o>0\text{ and }C_o\neq C',\\
0, & C_o=0.
\end{cases}
\end{equation}
The liquidation-incentive loss component is therefore
\begin{equation}
L_{\mathrm{liq}}=Q-R.
\end{equation}
It is combined with other scenario components such as direct collateral
impairment $C-C'$; the engine does not count the full debt principal as a
liquidation loss.
Other hard triggers include a live oracle breach, protocol emergency, or an
exit-liquidity threshold breach.  Hard triggers dominate cost comparison.

\section{Empirical 30-day probability}

For measurable price events, hourly observations are transformed into
overlapping 30-day windows of $720$ hours.  With $n$ valid windows and event
indicator $Y_k$, the point estimate is
\begin{equation}
\widehat p=\frac{k}{n},\qquad k=\sum_{k=1}^{n}Y_k.
\end{equation}
For example, a depeg event can be
\begin{equation}
Y_k=\ind\left[\min_{t\in W_k}p_t < 1-s\right],
\end{equation}
and a collateral shock can use peak-to-trough drawdown within $W_k$.

The 95\% Wilson interval, with $z=1.96$, is
\begin{align}
c&=\frac{\widehat p+z^2/(2n)}{1+z^2/n},\\
r&=\frac{z}{1+z^2/n}\sqrt{\frac{\widehat p(1-\widehat p)}{n}+\frac{z^2}{4n^2}},\\
[p_-,p_+]&=[\max(0,c-r),\min(1,c+r)].
\end{align}
Missing observations above 5\% of a window lower confidence; insufficiently
covered windows are excluded.  Overlapping windows are useful for the MVP but
are not independent, so Wilson intervals are descriptive rather than actuarial.

\section{Protocol evidence probability}

Rare smart-contract impairment cannot be estimated reliably from a small
incident count.  The MVP therefore computes a transparent evidence score
\begin{equation}
Q=0.30Q_{\mathrm{audit}}+0.25Q_{\mathrm{controls}}+0.20Q_{\mathrm{track}}
  +0.15Q_{\mathrm{dependencies}}+0.10Q_{\mathrm{monitoring}},
\end{equation}
with each component on $[0,100]$ and source provenance retained.  The score maps
to an explicitly model-estimated annual impairment probability
\begin{equation}
p_y(Q)=
\begin{cases}
0.005,&Q\geq85,\\
0.010,&70\leq Q<85,\\
0.030,&50\leq Q<70,\\
0.080,&Q<50.
\end{cases}
\end{equation}
The 30-day equivalent assumes a constant annual hazard:
\begin{equation}
p_{30}=1-(1-p_y)^{30/365}.
\end{equation}
Evidence-quality multipliers of 2, 3, or 4 widen uncertainty bands as evidence
weakens.  This result must be labelled \emph{model-estimated}; it is not a claim
that audit scores are observed failure frequencies.

\section{Exit-cost model}

Immediate exit cost is decomposed as
\begin{equation}
C_{\mathrm{exit}}=C_{\mathrm{fee}}+C_{\mathrm{slippage}}+C_{\mathrm{impact}}
 +C_{\mathrm{gas}}+C_{\mathrm{residual\ debt}}.
\end{equation}
For quoted input amount $q_{\mathrm{in}}$, reference mid-price $p$, and quoted
output $q_{\mathrm{out}}$, a transparent quote shortfall is
\begin{equation}
C_{\mathrm{quote}}=\max(0,q_{\mathrm{in}}p-q_{\mathrm{out}}),
\end{equation}
then split into explicit venue fee and residual price impact when the quote
adapter supplies both.  Gas cost in numeraire is
\begin{equation}
C_{\mathrm{gas}}=g\,b\,p_{\mathrm{native}},
\end{equation}
where $g$ is gas units, $b$ gas price in native token per unit, and
$p_{\mathrm{native}}$ the native-token numeraire price.  A quote older than 30
seconds or 20 blocks is refreshed once; if still expired, exit cost is unknown.

\section{Expected loss and hybrid recommendation}

For scenarios $j=1,\ldots,J$ with 30-day probabilities $p_j$ and impacts $L_j$,
the additive MVP expected loss is
\begin{equation}
\E[L_{30}]=\sum_{j=1}^{J}p_jL_j.
\end{equation}
This is a first-order approximation: correlated scenarios can double-count the
same tail event.  The radar graph and similarity matrices are intended to make
those common dependencies visible; a future model may replace the sum by a
joint scenario distribution.

The hybrid decision is evaluated in strict order:
\begin{equation}
D=
\begin{cases}
\texttt{UNAVAILABLE}, & \text{critical state, probability, or exit data missing/expired},\\
\texttt{EXIT}, & \exists j:T_j\neq\varnothing,\\
\texttt{UNAVAILABLE}, & \exists j:L_j>0\land p_j\text{ unavailable},\\
\texttt{EXIT}, & \E[L_{30}]>C_{\mathrm{exit}},\\
\texttt{STAY}, & \E[L_{30}]\leq C_{\mathrm{exit}}.
\end{cases}
\end{equation}
The explanation reports the decisive branch, scenario contributions, hard
triggers, exit-cost components, probability method, confidence, freshness, and
the affected graph path.  \texttt{UNAVAILABLE} is neutral and must never be
styled or described as safe.

\section{Known limitations and extension path}

\begin{itemize}
  \item The current similarity weights and health thresholds are policy values,
        not statistically calibrated universal constants.
  \item Peer-oracle comparison requires economic equivalence, unit normalization,
        and enough independent sources; otherwise the result is unknown.
  \item Overlapping historical windows induce dependence and rare-event samples
        are small.  Tail estimates require backtesting and stress calibration.
  \item Additive expected loss does not model correlated incidents or strategic
        reaction time.  A Bayesian network or Monte Carlo dependency model is a
        natural later extension.
  \item SQLite and in-process scheduling suit the local MVP.  Scaling to nine
        protocols and eight chains requires distributed ingestion, idempotent
        event processing, access control, durable notifications, and observability.
  \item The system is read-only decision support, not execution software and not
        investment advice.
\end{itemize}

\section{Implementation traceability}

The principal implementation locations are:
\begin{longtable}{>{\ttfamily}p{0.43\textwidth}p{0.49\textwidth}}
\toprule
services/risk/app/radar/domain & Similarity, oracle consensus, and health rules.\\
services/risk/app/radar/analysis & Incident lifecycle and reverse propagation.\\
services/risk/app/history/estimator.py & Empirical rolling windows and Wilson interval.\\
services/risk/app/evidence/scoring.py & Protocol evidence score and probability mapping.\\
services/risk/app/analysis & Scenario impact, exit cost, and hybrid recommendation.\\
services/risk/app/radar/scheduler.py & Fast/slow intervals and non-overlap policy.\\
apps/web/src/components & Institutional Terminal graph, exposure, incident, and notification UI.\\
\bottomrule
\end{longtable}

\end{document}
